% The twelve lines of the Hesse configuration, on the nine directions of
% Example ex:hesse.  There is no three-dimensional projection here and none is
% available: every one of the thirty-six pairs of directions lies on a line
% carrying a third, so a drawing by straight lines in the real plane would have
% no ordinary line, which Sylvester--Gallai forbids.  Two exact planar maps,
% both at 1.95cm per residue:
%   array   X =  1.95 p,   Y = -1.95 q       (p = slot, q = phase, each 0,1,2)
%   wrap    a step of the dashed class is (X,Y) -> (X + 1.95, Y - 1.95), of the
%           dotted class (X,Y) -> (X + 1.95, Y + 1.95); a step that leaves the
%           array is returned by Y -> Y + 5.85 resp. Y -> Y - 5.85, with
%           5.85 = 3 * 1.95, and is drawn as two stubs of 0.56 in X and in Y.
% Every coordinate below is 1.95 times its residue pair, written out with that
% pair in a trailing comment.  The dotted class is set at 0.75pt so that its
% round dots carry the weight of a 0.45pt stroke; the other three are 0.45pt.
\begin{tikzpicture}[
  x=1cm, y=1cm,
  gnode/.style ={circle, draw=black!70, fill=white, line width=0.5pt,
                 inner sep=0pt, minimum size=5.4mm, font=\footnotesize},
  plainline/.style={line width=0.45pt, black!65},
  dashline/.style ={line width=0.45pt, black!65,
                    dash pattern=on 1.4pt off 1.2pt},
  dotline/.style  ={line width=0.75pt, black!65, line cap=round,
                    dash pattern=on 0.1pt off 1.25pt},
  stub/.style  ={-{Latex[length=2mm,width=1.5mm]}},
  ring/.style  ={line width=0.5pt, black!55},
  tag/.style   ={fill=white, inner sep=1pt, font=\scriptsize},
  every node/.style={font=\small}]

% ---- the triple {u_0, u_1, u_3}, which lies on no line -------------------
\path[fill=black!12] ( 0.00, 0.00)         % u_0 = (p,q) = (0,0)
                  -- ( 0.00,-1.95)         % u_1 = (0,1)
                  -- ( 1.95, 0.00)         % u_3 = (1,0)
                  -- cycle;

% ---- the nine directions, column p and row q as in the ex:hesse display ---
\node[gnode] (u0) at ( 0.00, 0.00) {$u_0$};   % (p,q) = (0,0)
\node[gnode] (u1) at ( 0.00,-1.95) {$u_1$};   %         (0,1)
\node[gnode] (u2) at ( 0.00,-3.90) {$u_2$};   %         (0,2)
\node[gnode] (u3) at ( 1.95, 0.00) {$u_3$};   %         (1,0)
\node[gnode] (u4) at ( 1.95,-1.95) {$u_4$};   %         (1,1)
\node[gnode] (u5) at ( 1.95,-3.90) {$u_5$};   %         (1,2)
\node[gnode] (u6) at ( 3.90, 0.00) {$u_6$};   %         (2,0)
\node[gnode] (u7) at ( 3.90,-1.95) {$u_7$};   %         (2,1)
\node[gnode] (u8) at ( 3.90,-3.90) {$u_8$};   %         (2,2)

% ---- q constant: the three rows ------------------------------------------
\draw[plainline] (u0)--(u3) (u3)--(u6) (u1)--(u4) (u4)--(u7) (u2)--(u5) (u5)--(u8);
% ---- p constant: the three columns ---------------------------------------
\draw[plainline] (u0)--(u1) (u1)--(u2) (u3)--(u4) (u4)--(u5) (u6)--(u7) (u7)--(u8);

% ---- q - p constant: u_0u_4u_8, u_1u_5u_6, u_2u_3u_7 ---------------------
\draw[dashline] (u0)--(u4) (u4)--(u8) (u1)--(u5) (u3)--(u7);
\draw[dashline, stub] (u5)--( 2.51,-4.46);    % u_5 leaves, resumes at u_6
\draw[dashline, stub] ( 3.34, 0.56)--(u6);
\draw[dashline, stub] (u2)--( 0.56,-4.46);    % u_2 leaves, resumes at u_3
\draw[dashline, stub] ( 1.39, 0.56)--(u3);

% ---- q + p constant: u_2u_4u_6, u_0u_5u_7, u_1u_3u_8 ---------------------
\draw[dotline] (u2)--(u4) (u4)--(u6) (u5)--(u7) (u1)--(u3);
\draw[dotline, stub] (u0)--( 0.56, 0.56);     % u_0 leaves, resumes at u_5
\draw[dotline, stub] ( 1.39,-4.46)--(u5);
\draw[dotline, stub] (u3)--( 2.51, 0.56);     % u_3 leaves, resumes at u_8
\draw[dotline, stub] ( 3.34,-4.46)--(u8);

% ---- where each stub resumes ---------------------------------------------
\node[anchor=south, font=\scriptsize] at ( 0.56, 0.70) {$u_5$};
\node[anchor=south, font=\scriptsize] at ( 1.39, 0.70) {$u_2$};
\node[anchor=south, font=\scriptsize] at ( 2.51, 0.70) {$u_8$};
\node[anchor=south, font=\scriptsize] at ( 3.34, 0.70) {$u_5$};
\node[anchor=north, font=\scriptsize] at ( 0.56,-4.60) {$u_3$};
\node[anchor=north, font=\scriptsize] at ( 1.39,-4.60) {$u_0$};
\node[anchor=north, font=\scriptsize] at ( 2.51,-4.60) {$u_6$};
\node[anchor=north, font=\scriptsize] at ( 3.34,-4.60) {$u_3$};

% ---- the least eigenvalue, on one line of each of the four classes --------
\node[tag] at ( 0.98,-1.95) {$0$};            % on u_1u_4, q constant
\node[tag] at ( 3.90,-0.98) {$0$};            % on u_6u_7, p constant
\node[tag] at ( 2.50,-2.50) {$0$};            % on u_4u_8, q - p constant
\node[tag] at ( 0.49,-3.41) {$0$};            % on u_2u_4, q + p constant

% ---- the ringed triple ---------------------------------------------------
\draw[ring] (u0) circle (0.42);
\draw[ring] (u1) circle (0.42);
\draw[ring] (u3) circle (0.42);

% ---- the two residues ----------------------------------------------------
\node[anchor=south, font=\footnotesize] at ( 0.00, 1.14) {$p=0$};
\node[anchor=south, font=\footnotesize] at ( 1.95, 1.14) {$p=1$};
\node[anchor=south, font=\footnotesize] at ( 3.90, 1.14) {$p=2$};
\node[anchor=east, font=\footnotesize]  at (-0.72, 0.00) {$q=0$};
\node[anchor=east, font=\footnotesize]  at (-0.72,-1.95) {$q=1$};
\node[anchor=east, font=\footnotesize]  at (-0.72,-3.90) {$q=2$};

% ---- the four parallel classes -------------------------------------------
\draw[plainline] ( 4.70, 1.02)--( 5.20, 1.02);
\draw[plainline] ( 4.95, 0.64)--( 4.95, 0.24);
\draw[dashline]  ( 4.75, 0.06)--( 5.15,-0.34);
\draw[dotline]   ( 4.75,-0.92)--( 5.15,-0.52);
\node[anchor=west, font=\footnotesize] at ( 5.38, 1.02) {$q$ constant};
\node[anchor=west, font=\footnotesize] at ( 5.38, 0.44) {$p$ constant};
\node[anchor=west, font=\footnotesize] at ( 5.38,-0.14) {$q-p$ constant};
\node[anchor=west, font=\footnotesize] at ( 5.38,-0.72) {$q+p$ constant};

\node[anchor=north west, align=left, font=\scriptsize] at ( 4.70,-1.15)
  {$p$ slot, $q$ phase, modulo three\\[2pt]
   $3+3+3+3=12$ lines, four per point\\[2pt]
   $\binom{9}{3}=84=12+72$\\[5pt]
   the twelve, tagged $0$:\\[2pt]
   $\beta_{abc}=-\tfrac{27}{8}$, \ $\lmin(\SC)=0$\\[5pt]
   the seventy-two, $\{u_0,u_1,u_3\}$ ringed:\\[2pt]
   $\beta_{abc}=\tfrac{27}{16}(1\pm i\sqrt3)$,\\[2pt]
   $\lmin(\SC)=3\bigl(1-\cos\tfrac{2\pi}{9}\bigr)$\\[5pt]
   a stub marked $u_d$ resumes at $u_d$};

\end{tikzpicture}
